\documentclass[14pt, a4paper, landscape]{extarticle}

\usepackage{xcolor}
\usepackage{geometry}
\usepackage{graphicx}
\usepackage{booktabs}
\usepackage{array}
\usepackage{longtable}
\usepackage{titlesec}
\usepackage{enumitem}
\usepackage{fancyhdr}
\usepackage{tikz}

\geometry{margin=1cm}

\usepackage{xepersian}
\settextfont{Vazirmatn}
\setlatintextfont{Times New Roman}
\setdigitfont{Vazirmatn}

\definecolor{primary}{RGB}{139, 90, 43}
\definecolor{accent}{RGB}{224, 122, 63}
\definecolor{bg}{RGB}{245, 240, 232}
\definecolor{darkbg}{RGB}{92, 58, 26}

\pagestyle{fancy}
\fancyhf{}
\fancyhead[R]{\small\textcolor{primary}{مهندسی نرم‌افزار — جلسهٔ ۶}}
\fancyfoot[C]{\textcolor{primary}{\thepage}}
\renewcommand{\headrulewidth}{0.5pt}
\setlength{\headheight}{15pt}

\setlength{\parindent}{0pt}
\setlength{\parskip}{0.5em}

\newcommand{\slide}[1]{%
    \newpage
    \begin{center}
    \begin{minipage}[t]{0.95\textwidth}
    \vspace{0.5cm}
    #1
    \vspace{0.5cm}
    \end{minipage}
    \end{center}
}

\newcommand{\slidetitle}[1]{%
    \begin{center}
        {\LARGE\bfseries\textcolor{primary}{#1}}
    \end{center}
    \vspace{0.5cm}
}

\newcommand{\myaccent}[1]{\textcolor{accent}{#1}}

\begin{document}

\pagestyle{empty}

\begin{tikzpicture}[remember picture, overlay]
    \fill[darkbg] (current page.north west) rectangle (current page.south 
east);
\end{tikzpicture}

\begin{center}
    \vspace*{4cm}
    {\color{white}\Huge\bfseries مهندسی نرم‌افزار}
    
    \vspace{1cm}
    
    {\color{accent}\LARGE جلسهٔ ۶ — طراحی شیءگرا و SOLID}
    
    \vspace{2cm}
    
    {\color{white}\Large ابراهیم فضلی}
    
    \vspace{0.5cm}
    
    {\color{white}\large ترم ۱۴۰۴}
\end{center}

\pagestyle{fancy}

\slide{
    \slidetitle{مرور جلسهٔ ۵}
    \vspace{0.5cm}
    \begin{itemize}[itemsep=0.5em]
        \item Sprint Review: نمایش محصول، نه گزارش
        \item Definition of Done
        \item Retrospective: Start / Stop / Continue
        \item درس‌های Sprint 1
        \item برنامهٔ Sprint 2
    \end{itemize}
    \vspace{0.5cm}
    \begin{center}
        {\myaccent{امروز: چطور کد بهتر بنویسیم؟}}
    \end{center}
}

\slide{
    \slidetitle{چرا طراحی شیءگرا مهم است؟}
    \vspace{0.5cm}
    \begin{itemize}[itemsep=0.6em]
        \item {\large کد خوب = قابل نگهداری}
        \item {\large کد بد = هزینهٔ تغییر بالا}
        \item {\large طراحی شیءگرا = مدیریت پیچیدگی}
        \item {\large SOLID = اصول طراحی خوب}
    \end{itemize}
    \vspace{0.8cm}
    \begin{center}
        {\large\myaccent{۸۰٪ زمان توسعه، صرف خواندن کد می‌شود.}}
    \end{center}
}

\slide{
    \slidetitle{چهار اصل OOP}
    \vspace{0.5cm}
    \begin{center}
        \begin{tabular}{lp{10cm}}
            \toprule
            \textbf{اصل} & \textbf{معنی} \\
            \midrule
            Encapsulation & کپسوله‌سازی داده و رفتار \\
            \hline
            Inheritance & ارث‌بری از کلاس والد \\
            \hline
            Polymorphism & چندشکلی — یک نام، چند رفتار \\
            \hline
            Abstraction & انتزاع — پنهان کردن جزئیات \\
            \bottomrule
        \end{tabular}
    \end{center}
    \vspace{0.8cm}
    \begin{center}
        {\myaccent{این چهار اصل، پایهٔ طراحی شیءگرا هستند.}}
    \end{center}
}

\slide{
    \slidetitle{Encapsulation — کپسوله‌سازی}
    \vspace{0.5cm}
    {\large داده و متدها در یک واحد (کلاس) قرار می‌گیرند.}
    \vspace{0.8cm}
    \begin{itemize}[itemsep=0.5em]
        \item داده‌های داخلی \textbf{پنهان} هستند
        \item دسترسی فقط از طریق \textbf{متدهای عمومی}
        \item مزیت: تغییر داخلی، بدون تأثیر بیرونی
    \end{itemize}
    \vspace{0.8cm}
    \begin{center}
        {\myaccent{مثال: کلاس Product — قیمت خصوصی، متد getPrice() عمومی}}
    \end{center}
}

\slide{
    \slidetitle{Inheritance — ارث‌بری}
    \vspace{0.5cm}
    {\large کلاس فرزند، ویژگی‌های کلاس والد را به ارث می‌برد.}
    \vspace{0.8cm}
    \begin{itemize}[itemsep=0.5em]
        \item مزیت: \textbf{استفادهٔ مجدد} از کد
        \item عیب: \textbf{وابستگی} زیاد
        \item توصیه: \textbf{ترکیب} بهتر از ارث‌بری است
    \end{itemize}
    \vspace{0.8cm}
    \begin{center}
        {\myaccent{مثال: Admin از User ارث می‌برد}}
    \end{center}
}

\slide{
    \slidetitle{Polymorphism — چندشکلی}
    \vspace{0.5cm}
    {\large یک نام، چند رفتار.}
    \vspace{0.8cm}
    \begin{itemize}[itemsep=0.5em]
        \item متد \textbf{مشترک}، پیاده‌سازی \textbf{متفاوت}
        \item مزیت: کد \textbf{عمومی}، رفتار \textbf{خاص}
        \item مثال: \lr{pay()} برای کارت، PayPal، درگاه
    \end{itemize}
    \vspace{0.8cm}
    \begin{center}
        {\myaccent{یک interface، چند implementation}}
    \end{center}
}

\slide{
    \slidetitle{Abstraction — انتزاع}
    \vspace{0.5cm}
    {\large پنهان کردن جزئیات پیاده‌سازی.}
    \vspace{0.8cm}
    \begin{itemize}[itemsep=0.5em]
        \item کاربر \textbf{چه کاری} انجام می‌دهد را می‌بیند
        \item \textbf{چطور} انجام می‌دهد را نمی‌بیند
        \item ابزار: Interface، Abstract Class
    \end{itemize}
    \vspace{0.8cm}
    \begin{center}
        {\myaccent{مثال: UserRepository.save() — بدون دانستن جزئیات 
دیتابیس}}
    \end{center}
}

\slide{
    \slidetitle{SOLID چیست؟}
    \vspace{0.5cm}
    \begin{center}
        {\large\bfseries پنج اصل طراحی شیءگرا}
    \end{center}
    \vspace{0.8cm}
    \begin{center}
        \begin{tabular}{lp{9cm}}
            \toprule
            \textbf{حرف} & \textbf{اصل} \\
            \midrule
            \textbf{S} & Single Responsibility Principle \\
            \hline
            \textbf{O} & Open/Closed Principle \\
            \hline
            \textbf{L} & Liskov Substitution Principle \\
            \hline
            \textbf{I} & Interface Segregation Principle \\
            \hline
            \textbf{D} & Dependency Inversion Principle \\
            \bottomrule
        \end{tabular}
    \end{center}
}

\slide{
    \slidetitle{S — Single Responsibility}
    \vspace{0.5cm}
    \begin{center}
        {\large\bfseries هر کلاس فقط یک مسئولیت}
    \end{center}
    \vspace{0.8cm}
    \begin{itemize}[itemsep=0.5em]
        \item کلاس \textbf{فقط یک دلیل} برای تغییر داشته باشد
        \item اگر کلاس هم محاسبه، هم ذخیره، هم چاپ کند → نقض SRP
    \end{itemize}
    \vspace{0.8cm}
    \begin{center}
        {\myaccent{مثال: Order نباید فاکتور چاپ کند. Invoice جدا باشد.}}
    \end{center}
}

\slide{
    \slidetitle{O — Open/Closed}
    \vspace{0.5cm}
    \begin{center}
        {\large\bfseries باز برای توسعه، بسته برای تغییر}
    \end{center}
    \vspace{0.8cm}
    \begin{itemize}[itemsep=0.5em]
        \item افزودن ویژگی \textbf{جدید} بدون تغییر کد \textbf{قدیمی}
        \item ابزار: Interface، Abstract Class
    \end{itemize}
    \vspace{0.8cm}
    \begin{center}
        {\myaccent{مثال: افزودن روش پرداخت جدید، بدون تغییر کد اصلی}}
    \end{center}
}

\slide{
    \slidetitle{L — Liskov Substitution}
    \vspace{0.5cm}
    \begin{center}
        {\large\bfseries فرزند باید جایگزین والد شود}
    \end{center}
    \vspace{0.8cm}
    \begin{itemize}[itemsep=0.5em]
        \item هر جا \lr{User} انتظار می‌رود، \lr{Admin} هم کار کند
        \item فرزند نباید \textbf{رفتار} والد را نقض کند
    \end{itemize}
    \vspace{0.8cm}
    \begin{center}
        {\myaccent{مثال بد: Square از Rectangle ارث می‌برد، ولی setWidth() 
را نقض می‌کند}}
    \end{center}
}

\slide{
    \slidetitle{I — Interface Segregation}
    \vspace{0.5cm}
    \begin{center}
        {\large\bfseries رابط‌های کوچک، نه بزرگ}
    \end{center}
    \vspace{0.8cm}
    \begin{itemize}[itemsep=0.5em]
        \item کلاس نباید متدهای \textbf{غیرضروری} را پیاده کند
        \item چند interface کوچک بهتر از یک interface بزرگ
    \end{itemize}
    \vspace{0.8cm}
    \begin{center}
        {\myaccent{مثال: PaymentGateway و RefundGateway جدا باشند}}
    \end{center}
}

\slide{
    \slidetitle{D — Dependency Inversion}
    \vspace{0.5cm}
    \begin{center}
        {\large\bfseries به انتزاع وابسته باش، نه پیاده‌سازی}
    \end{center}
    \vspace{0.8cm}
    \begin{itemize}[itemsep=0.5em]
        \item کلاس سطح بالا نباید به کلاس سطح پایین وابسته باشد
        \item هر دو به \textbf{interface} وابسته باشند
    \end{itemize}
    \vspace{0.8cm}
    \begin{center}
        {\myaccent{مثال: OrderService به \lr{IOrderRepository} وابسته 
باشد، نه به \lr{PostgresOrderRepository}}}
    \end{center}
}

\slide{
    \slidetitle{مثال: نقض SRP در HomeStore}
    \vspace{0.5cm}
    {\large کلاس \lr{Order} که همه‌کاره است:}
    \vspace{0.5cm}
    \begin{itemize}[itemsep=0.4em]
        \item محاسبهٔ مجموع
        \item ذخیره در دیتابیس
        \item ارسال ایمیل تأیید
        \item چاپ فاکتور
        \item اطلاع به انبار
    \end{itemize}
    \vspace{0.8cm}
    \begin{center}
        {\large\myaccent{۵ دلیل برای تغییر = نقض شدید SRP}}
    \end{center}
}

\slide{
    \slidetitle{Refactor با SOLID}
    \vspace{0.5cm}
    {\large جدا کردن مسئولیت‌ها:}
    \vspace{0.5cm}
    \begin{center}
        \begin{tabular}{ll}
            \toprule
            \textbf{کلاس} & \textbf{مسئولیت} \\
            \midrule
            \lr{Order} & محاسبهٔ مجموع \\
            \hline
            \lr{OrderRepository} & ذخیره در دیتابیس \\
            \hline
            \lr{EmailService} & ارسال ایمیل \\
            \hline
            \lr{InvoiceService} & چاپ فاکتور \\
            \hline
            \lr{StockService} & اطلاع به انبار \\
            \bottomrule
        \end{tabular}
    \end{center}
    \vspace{0.5cm}
    \begin{center}
        {\myaccent{هر کلاس = یک مسئولیت}}
    \end{center}
}

\slide{
    \slidetitle{Coupling و Cohesion}
    \vspace{0.5cm}
    \begin{center}
        \begin{tabular}{lp{9cm}}
            \toprule
            \textbf{مفهوم} & \textbf{توضیح} \\
            \midrule
            Coupling & وابستگی بین کلاس‌ها — \textbf{کم} بهتر \\
            \hline
            Cohesion & ارتباط داخلی کلاس — \textbf{زیاد} بهتر \\
            \bottomrule
        \end{tabular}
    \end{center}
    \vspace{0.8cm}
    \begin{itemize}[itemsep=0.4em]
        \item \textbf{Coupling کم:} تغییر یک کلاس، بقیه را نشکند
        \item \textbf{Cohesion زیاد:} همهٔ متدهای کلاس، یک هدف
    \end{itemize}
}

\slide{
    \slidetitle{مثال: Coupling در HomeStore}
    \vspace{0.5cm}
    {\large کد با Coupling بالا:}
    \vspace{0.5cm}
    \begin{itemize}[itemsep=0.4em]
        \item \lr{OrderService} مستقیماً \lr{PostgresOrderRepository} را 
می‌سازد
        \item تغییر دیتابیس = تغییر OrderService
    \end{itemize}
    \vspace{0.8cm}
    {\large کد با Coupling پایین:}
    \vspace{0.5cm}
    \begin{itemize}[itemsep=0.4em]
        \item \lr{OrderService} به \lr{IOrderRepository} وابسته است
        \item تغییر دیتابیس = فقط تغییر implementation
    \end{itemize}
}

\slide{
    \slidetitle{Design Patterns — مقدمه}
    \vspace{0.5cm}
    \begin{center}
        {\large\bfseries راه‌حل‌های آماده برای مشکلات رایج}
    \end{center}
    \vspace{0.8cm}
    \begin{itemize}[itemsep=0.5em]
        \item \textbf{Creational:} Factory، Singleton، Builder
        \item \textbf{Structural:} Adapter، Decorator، Facade
        \item \textbf{Behavioral:} Strategy، Observer، Command
    \end{itemize}
    \vspace{0.8cm}
    \begin{center}
        {\myaccent{در جلسهٔ ۷ عمیق‌تر می‌بینیم.}}
    \end{center}
}

\slide{
    \slidetitle{تمرین: Refactor کد}
    \vspace{0.5cm}
    {\large کد زیر را با SOLID بازنویسی کنید:}
    \vspace{0.5cm}
    \begin{center}
        \begin{tabular}{lp{9cm}}
            \toprule
            \textbf{کلاس} & \textbf{مسئولیت‌ها} \\
            \midrule
            \lr{ProductManager} & افزودن، ویرایش، حذف، جست‌وجو، گزارش \\
            \hline
            \lr{OrderManager} & ثبت، لغو، پیگیری، ایمیل، فاکتور \\
            \bottomrule
        \end{tabular}
    \end{center}
    \vspace{0.5cm}
    {\large \textbf{سؤال:} چند کلاس جدید لازم است؟ هر کلاس چه مسئولیتی 
دارد؟}
}

\slide{
    \slidetitle{تسک هفتهٔ آینده}
    \vspace{0.5cm}
    \begin{minipage}[t]{0.48\textwidth}
        {\bfseries\textcolor{primary}{تسک شما:}}
        \vspace{0.3cm}
        \begin{itemize}[itemsep=0.3em]
            \item کد Sprint 1 را بازبینی کنید
            \item کلاس‌های بزرگ را جدا کنید
            \item SRP را رعایت کنید
            \item Coupling را کم کنید
            \item مستندسازی تصمیم‌ها
        \end{itemize}
    \end{minipage}
    \hfill
    \begin{minipage}[t]{0.48\textwidth}
        {\bfseries\myaccent{تسک من:}}
        \vspace{0.3cm}
        \begin{itemize}[itemsep=0.3em]
            \item Refactor کد HomeStore
            \item مستند SOLID
            \item جزوهٔ مطالعه
            \item آماده‌سازی جلسهٔ ۷
        \end{itemize}
    \end{minipage}
}

\slide{
    \slidetitle{سه برداشت کلیدی}
    \vspace{1cm}
    \begin{enumerate}[itemsep=0.8em]
        \item {\large SOLID = پنج اصل برای کد قابل نگهداری}
        \item {\large Coupling کم، Cohesion زیاد}
        \item {\large Refactor یک عادت است، نه یک پروژه}
    \end{enumerate}
}

\slide{
    \begin{center}
        \vspace{3cm}
        {\Huge\bfseries سؤالی دارید؟}
        \vspace{1.5cm}
        {\LARGE\myaccent{به امید یک ترم پربار}}
    \end{center}
}

\end{document}
